% ----- Event of a timeline -----
% \event[dx]{year}{citation keys}{text}{text width}{min height}{y}{font}{fill}
% The optional dx (default 0) shifts the box horizontally from its year, so that
% boxes of close years do not overlap; the dashed line stays on the year.
\newcommand{\event}[9][0]{%
    \pgfmathsetmacro{\x}{(#2 - \timelineStart) * \timelineScale}%
    \pgfmathsetmacro{\y}{#7}%
    \begin{scope}[on background layer]
        \draw[thick, dashed] (\x,0) -- ++(0,\y);
    \end{scope}
    \node[
        draw, rounded corners, very thick, fill=#9,
        minimum height=#6, text width=#5, align=center, font=#8,
        execute at begin node={\hyphenpenalty=10000\exhyphenpenalty=10000}
    ] at (\x+#1,\y) {\textbf{\textcite{#3}} \\[2pt] #4};
}

% ----- Main timeline command -----
% Usage: \Timeline[start_year,end_year,scale]{ list of \event{...} ... }
\newcommand{\Timeline}[2][1943,2025,0.3]{%
    % Extract parameters
    \edef\timelineParams{#1}%
    \pgfmathsetmacro{\timelineStart}{{\timelineParams}[0]}%
    \pgfmathsetmacro{\timelineEnd}{{\timelineParams}[1]}%
    \pgfmathsetmacro{\timelineScale}{{\timelineParams}[2]}%
    %
    \begin{tikzpicture}
        % Main timeline arrow
        \draw[draw, thick, -{Stealth[length=3mm]}, line cap=round] 
            (0,0) -- ({(\timelineEnd - \timelineStart) * \timelineScale},0);
        % All events
        #2
    \end{tikzpicture}%
}